\chapter{创新性说明}\label{ux7b2cux4e94ux7ae0-ux521bux65b0ux6027ux8bf4ux660e}

前文从系统设计与实现角度介绍了V-Guard的整体架构与核心功能，并在第四章通过实验验证了其在语义保护和身份保护两个维度上的有效性。在上述设计与实验结果的基础上，本章进一步从技术创新的角度，对本作品的核心创新点进行专门说明：

\textbf{（1）面向声音资产发布前的源头主动防护。}语音克隆攻击的典型流程中，攻击者从公开渠道获取语音素材后，通常依赖ASR系统自动生成文本标注，再通过``音频-文本''对微调预训练模型完成克隆。这一流程不依赖人工标注，也不依赖于攻破检测系统，可以在事后检测部署之前就完成样本采集和模型适配。事后检测方法只能在克隆完成后进行识别和追责，对攻击前期的数据收集和模型训练阶段没有约束力。因此，当防护措施停留在检测层面时，攻击者总可以在数据源头绕过监管。

针对上述问题，V-Guard将防护前移至语音公开发布环节。用户发布语音之前，系统对原始音频施加保护扰动，使音频在语义表示和身份表示空间上偏离原始样本，但人耳听感基本不变。攻击者即便采集到发布后的音频，也难以从中提取稳定的语义内容和说话人身份信息来支撑克隆建模。源头防护不依赖于对攻击者所使用的克隆模型或ASR类型的先验知识，能够从样本层面覆盖端到端克隆链条的起点，在数据被利用之前降低其可用价值，与事后检测形成主动与被动两种不同的防御逻辑。

\textbf{（2）面向语义表示与声音身份表示的联合保护。}语音克隆系统对参考音频的利用涉及语义内容的编码和说话人身份的声学特征编码两个层面。已有防护工作通常只针对其中一个层面进行干扰，当攻击所需的两类信息中只有一类被干扰时，攻击者仍可利用未被保护的那类信息完成克隆。

本工作将语义表示和声音身份表示纳入同一框架，使同一个保护扰动同时作用于两个表示空间。在语义链路上，防护目标从ASR的输出文本层面下沉至语音编码器的中间表示层，利用离散语音token量化过程对连续表示变化的敏感性，使扰动引发token序列的结构性改变，进而影响下游语音大模型对内容的整体理解，这种扰动不依赖于特定ASR的解码策略，在不同识别系统之间具备更强的迁移性。在身份链路上，采用VITS、GPT-SoVITS、CosyVoice
CAM++、StyleTTS2风格编码器、WavLM和MFCC等六类编码器作为代理模型，覆盖变分自编码器、扩散模型、流匹配和传统声学特征等主流架构，使扰动同时在多个身份表示空间中偏离原始说话人，避免对单一编码器过拟合，使攻击者无论使用哪一种编码器进行音色条件提取，均会遭遇明显的身份信息退化。此外，系统支持无目标和有目标两种保护模式，前者使克隆结果与原始说话人无关，后者将身份引导至特定目标方向，以适应不同防护需求。

两条链路共享同一个保护扰动，在优化过程中同时考虑语义表示偏移和身份表示偏移两个方向，使单一扰动能够同时破坏克隆语音的内容可理解性和身份可识别性。

\textbf{（3）心理声学约束下的高保真扰动生成。}本工作在扰动优化过程中引入心理声学掩蔽模型，实时计算原始语音的频域掩蔽阈值，将扰动能量约束在该阈值之下，使其主要分布于人耳听觉不敏感的区域。相比于单纯依赖幅度约束的扰动方法，心理声学掩蔽利用了听觉感知的结构性特点，能够在相同扰动幅度下获得更低的感知显著性。扰动优化目标中同时包含面向机器模型的语义和身份损失、面向人耳感知的心理声学损失以及~L2\emph{L}2\hspace{0pt}~幅度约束，使生成过程在防护强度与音频质量之间形成可量化的平衡。受保护音频在信噪比和语音质量感知评估等客观指标上保持在可用范围之内，同时仍能对下游克隆模型产生有效的语义和身份干扰。

\textbf{（4）防护-评估-导出全流程系统架构。}本系统将主动防护算法、下游攻击模拟和多维评估能力整合为统一的网页端闭环流程。用户上传音频后，系统自动完成格式预处理、保护扰动生成、ASR转写对比、说话人相似度评估和克隆防护验证，并同步展示原始与保护音频在语义识别和身份提取两个维度上的差异。该闭环设计使防护效果可在同一系统内完成验证，无需额外搭建评估环境或依赖人工主观判断；历史任务管理和报告导出功能则为防护结果的追溯、复核和归档提供了支撑。这一架构将主动防护从单一的扰动生成扩展为可验证、可追溯的完整流程。

\FloatBarrier
